\subsection{Augmented Assignment}

Augmented assignment is assignment combined with an operation.

Instead of writing:

\begin{verbatim}
score = score + 10
\end{verbatim}

Python allows:

\begin{verbatim}
score += 10
\end{verbatim}

The second form is shorter, but it should not be misunderstood. For immutable singular values such as integers, floats, booleans, strings, and complex numbers, augmented assignment does not edit the object in place. It reads the old object, computes a new result object, and rebinds the name.

The basic forms are:

\begin{verbatim}
x += y
x -= y
x *= y
x /= y
x //= y
x %= y
x **= y
\end{verbatim}

For now, the important conceptual form is:

\begin{verbatim}
name = name operator value
\end{verbatim}

So:

\begin{verbatim}
count += 1
\end{verbatim}

means approximately:

\begin{verbatim}
count = count + 1
\end{verbatim}

The approximation is good enough for immutable singular values. Later, mutable containers will make the story more delicate.

\subsubsection{Numeric Augmented Assignment (\texttt{+=}, \texttt{-=}, \texttt{*=}, \texttt{/=}) as Read-Operate-Rebind for Immutable Singular Values}

Start with a plain integer.

\begin{verbatim}
score = 42

print("before")
print(f"score: {score}")
print(f"type:  {type(score)}")
print(f"id:    {id(score)}")

old_id = id(score)

score += 8

print()
print("after score += 8")
print(f"score: {score}")
print(f"type:  {type(score)}")
print(f"id:    {id(score)}")
print(f"same object as before: {id(score) == old_id}")
\end{verbatim}

Possible output:

\begin{verbatim}
before
score: 42
type:  <class 'int'>
id:    140734432946520

after score += 8
score: 50
type:  <class 'int'>
id:    140734432946776
same object as before: False
\end{verbatim}

The exact \texttt{id()} values are not important. The important point is this: the integer object \texttt{42} was not edited into \texttt{50}. Integers are immutable. Python computed the result of \texttt{42 + 8}, then rebound the name \texttt{score} to the resulting integer object.

Conceptually:

\begin{verbatim}
before:
    score ---> int object 42

operation:
    42 + 8 produces int object 50

after:
    score ---> int object 50
\end{verbatim}

The integer object has methods and special methods that support this behavior.

\begin{verbatim}
value = 42

print(f"type(value): {type(value)}")

print("selected int methods:")
print(f"{'bit_length':<15} {'bit_length' in dir(value)}")
print(f"{'__add__':<15} {'__add__' in dir(value)}")
print(f"{'__sub__':<15} {'__sub__' in dir(value)}")
print(f"{'__mul__':<15} {'__mul__' in dir(value)}")
print(f"{'__truediv__':<15} {'__truediv__' in dir(value)}")

print()
print(f"value.bit_length(): {value.bit_length()}")
print(f"value + 8:          {value + 8}")
print(f"value * 2:          {value * 2}")
\end{verbatim}

Possible output:

\begin{verbatim}
type(value): <class 'int'>
selected int methods:
bit_length      True
__add__         True
__sub__         True
__mul__         True
__truediv__     True

value.bit_length(): 6
value + 8:          50
value * 2:          84
\end{verbatim}

The method \texttt{bit\_length()} is an integer-specific method. The double-underscore methods such as \texttt{\_\_add\_\_}, \texttt{\_\_sub\_\_}, and \texttt{\_\_mul\_\_} are part of the object protocol that connects operators to object behavior. In normal code, we write \texttt{value + 8}, not \texttt{value.\_\_add\_\_(8)}.

Several numeric augmented assignments:

\begin{verbatim}
number = 42

print(f"start:      {number}")

number += 10
print(f"after +=:   {number}")

number -= 2
print(f"after -=:   {number}")

number *= 3
print(f"after *=:   {number}")

number /= 5
print(f"after /=:   {number}")

print()
print(f"final type: {type(number)}")
\end{verbatim}

Possible output:

\begin{verbatim}
start:      42
after +=:   52
after -=:   50
after *=:   150
after /=:   30.0

final type: <class 'float'>
\end{verbatim}

The last result is important. The operator \texttt{/} performs true division. Even when the mathematical result is whole, Python returns a floating-point object.

So this:

\begin{verbatim}
number /= 5
\end{verbatim}

does not mean: keep \texttt{number} as an integer and divide its storage by five.

It means:

\begin{verbatim}
number = number / 5
\end{verbatim}

If \texttt{number / 5} produces a \texttt{float}, then the name \texttt{number} is rebound to a \texttt{float} object.

A direct demonstration:

\begin{verbatim}
value = 10

print("before")
print(f"value: {value!r}")
print(f"type:  {type(value)}")

value /= 2

print()
print("after value /= 2")
print(f"value: {value!r}")
print(f"type:  {type(value)}")
\end{verbatim}

Possible output:

\begin{verbatim}
before
value: 10
type:  <class 'int'>

after value /= 2
value: 5.0
type:  <class 'float'>
\end{verbatim}

The name stayed the same. The referenced object changed from an integer object to a floating-point object.

For integer-preserving division, floor division uses \texttt{//}:

\begin{verbatim}
value = 10

value //= 2

print(f"value: {value!r}")
print(f"type:  {type(value)}")
\end{verbatim}

Possible output:

\begin{verbatim}
value: 5
type:  <class 'int'>
\end{verbatim}

But \texttt{//} means floor division, not exact mathematical division. With values that do not divide evenly, the result is rounded down toward negative infinity.

\begin{verbatim}
value = 11

value //= 2

print(f"value: {value}")
print(f"type:  {type(value)}")
\end{verbatim}

Possible output:

\begin{verbatim}
value: 5
type:  <class 'int'>
\end{verbatim}

Floating-point values behave similarly in the assignment model. The old object is not edited. A new numeric result is produced, and the name is rebound.

\begin{verbatim}
price = 19.99

print("before")
print(f"price: {price}")
print(f"type:  {type(price)}")
print(f"id:    {id(price)}")

old_id = id(price)

price *= 2

print()
print("after price *= 2")
print(f"price: {price}")
print(f"type:  {type(price)}")
print(f"id:    {id(price)}")
print(f"same object as before: {id(price) == old_id}")
\end{verbatim}

Possible output:

\begin{verbatim}
before
price: 19.99
type:  <class 'float'>
id:    2144421961872

after price *= 2
price: 39.98
type:  <class 'float'>
id:    2144421962416
same object as before: False
\end{verbatim}

A float object also exposes operator-related methods.

\begin{verbatim}
price = 19.99

print(f"type(price): {type(price)}")

print("selected float methods:")
print(f"{'as_integer_ratio':<20} {'as_integer_ratio' in dir(price)}")
print(f"{'is_integer':<20} {'is_integer' in dir(price)}")
print(f"{'__add__':<20} {'__add__' in dir(price)}")
print(f"{'__mul__':<20} {'__mul__' in dir(price)}")

print()
print(f"price.is_integer(): {price.is_integer()}")
print(f"price.as_integer_ratio(): {price.as_integer_ratio()}")
\end{verbatim}

Possible output:

\begin{verbatim}
type(price): <class 'float'>
selected float methods:
as_integer_ratio     True
is_integer           True
__add__              True
__mul__              True

price.is_integer(): False
price.as_integer_ratio(): (5626684784442573, 281474976710656)
\end{verbatim}

The strange-looking ratio is a reminder that many decimal-looking floating-point values are stored as binary approximations. This will matter more in the later floating-point subsection.

Augmented assignment also works with strings, but strings are immutable too.

\begin{verbatim}
line = "D'oh"

print("before")
print(f"line: {line!r}")
print(f"type: {type(line)}")
print(f"id:   {id(line)}")

old_id = id(line)

line += "!"

print()
print("after line += '!'")
print(f"line: {line!r}")
print(f"type: {type(line)}")
print(f"id:   {id(line)}")
print(f"same object as before: {id(line) == old_id}")
\end{verbatim}

Possible output:

\begin{verbatim}
before
line: "D'oh"
type: <class 'str'>
id:   2144421959120

after line += '!'
line: "D'oh!"
type: <class 'str'>
id:   2144421959280
same object as before: False
\end{verbatim}

The original string object was not edited. A new string object was produced.

For immutable singular values, the safe rule is:

\begin{quote}
Augmented assignment reads the current object, performs the operation, and rebinds the name to the result.
\end{quote}

\subsubsection{Mutation vs. Rebinding Preview for Later Mutable Containers}

The previous rule is reliable for immutable singular values. However, Python also has mutable objects. A mutable object can be changed in place.

Lists are mutable containers. They will be studied properly later, but they are useful here as a preview because they show why augmented assignment must be understood carefully.

First, inspect a list object:

\begin{verbatim}
scores = [10, 20, 30]

print(f"scores: {scores}")
print(f"type(scores): {type(scores)}")

print("selected list methods:")
print(f"{'append':<12} {'append' in dir(scores)}")
print(f"{'extend':<12} {'extend' in dir(scores)}")
print(f"{'pop':<12} {'pop' in dir(scores)}")
print(f"{'__iadd__':<12} {'__iadd__' in dir(scores)}")
\end{verbatim}

Possible output:

\begin{verbatim}
scores: [10, 20, 30]
type(scores): <class 'list'>
selected list methods:
append       True
extend       True
pop          True
__iadd__     True
\end{verbatim}

The method \texttt{append()} adds one element to the existing list. The method \texttt{extend()} adds several elements. The special method \texttt{\_\_iadd\_\_} is connected to in-place addition behavior for lists.

Now compare rebinding with mutation.

\begin{verbatim}
numbers = [10, 20]
alias = numbers

print("before")
print(f"numbers: {numbers}")
print(f"alias:   {alias}")
print(f"id(numbers): {id(numbers)}")
print(f"id(alias):   {id(alias)}")
print(f"numbers is alias: {numbers is alias}")

numbers += [30]

print()
print("after numbers += [30]")
print(f"numbers: {numbers}")
print(f"alias:   {alias}")
print(f"id(numbers): {id(numbers)}")
print(f"id(alias):   {id(alias)}")
print(f"numbers is alias: {numbers is alias}")
\end{verbatim}

Possible output:

\begin{verbatim}
before
numbers: [10, 20]
alias:   [10, 20]
id(numbers): 2144421972096
id(alias):   2144421972096
numbers is alias: True

after numbers += [30]
numbers: [10, 20, 30]
alias:   [10, 20, 30]
id(numbers): 2144421972096
id(alias):   2144421972096
numbers is alias: True
\end{verbatim}

Here, the list object itself changed. Both names still refer to the same list object, and both names see the changed content.

Conceptually:

\begin{verbatim}
before:
    numbers ----\
                 ---> list object [10, 20]
    alias   ----/

after numbers += [30]:
    numbers ----\
                 ---> same list object [10, 20, 30]
    alias   ----/
\end{verbatim}

This is mutation, not simple rebinding.

Now compare with ordinary \texttt{+} followed by assignment:

\begin{verbatim}
numbers = [10, 20]
alias = numbers

print("before")
print(f"numbers: {numbers}")
print(f"alias:   {alias}")
print(f"id(numbers): {id(numbers)}")
print(f"id(alias):   {id(alias)}")

numbers = numbers + [30]

print()
print("after numbers = numbers + [30]")
print(f"numbers: {numbers}")
print(f"alias:   {alias}")
print(f"id(numbers): {id(numbers)}")
print(f"id(alias):   {id(alias)}")
print(f"numbers is alias: {numbers is alias}")
\end{verbatim}

Possible output:

\begin{verbatim}
before
numbers: [10, 20]
alias:   [10, 20]
id(numbers): 2144421972096
id(alias):   2144421972096

after numbers = numbers + [30]
numbers: [10, 20, 30]
alias:   [10, 20]
id(numbers): 2144421972288
id(alias):   2144421972096
numbers is alias: False
\end{verbatim}

This time, the expression:

\begin{verbatim}
numbers + [30]
\end{verbatim}

created a new list object. Then the name \texttt{numbers} was rebound to that new object. The name \texttt{alias} still points to the original list.

Conceptually:

\begin{verbatim}
before:
    numbers ----\
                 ---> list object [10, 20]
    alias   ----/

after numbers = numbers + [30]:
    numbers ---> new list object [10, 20, 30]

    alias   ---> old list object [10, 20]
\end{verbatim}

So for mutable containers, these two forms can differ:

\begin{verbatim}
numbers += [30]
numbers = numbers + [30]
\end{verbatim}

For immutable singular values, they usually behave conceptually the same: compute a result and rebind the name.

For mutable containers, augmented assignment may mutate the existing object in place.

This is only a preview. The deeper rule will be developed later when lists, dictionaries, sets, and shared references are studied properly.

For now, keep this practical distinction:

\begin{itemize}
    \item \texttt{x = x + y} normally computes a result object and binds \texttt{x} to that result.
    \item \texttt{x += y} asks the object whether it can update itself in place.
    \item immutable objects cannot be updated in place, so the name is rebound to a result object;
    \item mutable objects may update themselves in place, so other names sharing the same object may see the change.
\end{itemize}

The beginner-safe rule is:

\begin{quote}
For numbers and strings, read augmented assignment as read-operate-rebind. For mutable containers, do not assume this; check whether the operation mutates the object in place.
\end{quote}